\scanpage{12}
\mspage{\handwritten{26}}

\uline{Théorème} \deleted{\unsure{3}}\handwritten{4}.6. Soient $\mathrm A$ un anneau préadmissible muni d'un idéal de définition $\mathrm I=\mathrm I_{\unsure{0|o}}$ et d'un système fondamental $(\mathrm I_i)_i$ d'idéaux d e définition $\mathrm I_0$, $\mathrm B$ une $\mathrm A$-algèbre topologique. On suppose\deleted{\unsure{r}} les $\mathrm A_i=\mathrm A/\mathrm I_i\mathrm A$ noethériens, et de plus les conditions suivantes sur $\mathrm B$:
\deleted{\unsure{a) Les $\mathrm A_i=\mathrm A/\mathrm I_i$ noethériens|a) les $\mathrm A_i$ sont nœthériens}}

\deleted{\unsure{b}}\handwritten{a}) Les \deleted{\unsure{B|M}}\handwritten{m}odule\deleted{s} topologique\deleted{s} \deleted{$\unsure{\mathrm B_i=\mathrm B/\mathrm I_i\mathrm B}$} \deleted{\unsure{sur} $\unsure{\mathrm A_i}$} \inserted{$\mathrm B$ sur $\mathrm A$} satisf\deleted{on}\handwritten{ai}t la condition $\mathrm P$.

\inserted{\handwritten{b) $B_0$ est formellement simple sur $A_0$}}

\deleted{b}\handwritten{c}) $\mathrm B$ est adique \deleted{pour un idéal contenant $\mathrm{IB}$}

\handwritten{d}) Les idéaux $\mathrm I_i\mathrm B$ dans $\mathrm B$ sont fermés (\deleted{inutile}\inserted{\handwritten{\unsure{automatiquement} vérifié}} si $\mathrm B$ est noethérien).

Sous ces conditions, on a ce qui suit\deleted{\unsure{e}}:

(i) $\mathrm B$ est formellement simple sur $\mathrm A$.

(ii) Soit $\mathrm C$ une $\mathrm A$-algèbre topologique satisfaisant à \deleted{\unsure{b}}\handwritten{c}) et \handwritten{d}) (mais pas nécessairement à a) ), et soit $u_0:\mathrm B_0\handwritten{\overset{\sim}{\to}}\mathrm C_0$ un isomorphisme d'al\`{g}èbres topologiques. Alors il existe un $\mathrm A$-homomorphisme continu $u:\mathrm B\handwritten{\to}\mathrm C$ se réduisant suivant $u_0$ \handwritten{(}en vertu de (i) et de 1.7.). Ce dernier étant choisi, les conditions suivantes sont équivalentes: \deleted{\unsure{ce}}

\deleted{\unsure{a)}}1) $\mathrm B$ est formellement simple sur $\mathrm A$.

2) \deleted{$\mathrm B$ est} $\mathrm B$ satisfait la c\deleted{\unsure{o}}ondition a) ci-dessus, i.e. \deleted{est} satisf\deleted{\unsure{i|.}}ait la condition $\mathrm P$.

3) Les $\mathrm B_i$ sont plats sur $\mathrm A_i$.

4) L'homomorphisme $u:\mathrm B\handwritten{\to}\mathrm C$ est un isomorphisme\deleted{\unsure{,}} de $\mathrm A$-algèbres top.

\inserted{\handwritten{5) $\deleted{\unsure{B}}C$ est isom. à $B$ comme $A$-algèbre top.}}

Démonstration. L'assertion (i) résulte déjà de a)\inserted{\handwritten{ et b)}}, grâce à 3.1.

Prouvons donc (ii). En vertu de (i) d\deleted{\unsure{q}}éjà prouvé, on a 2)$\handwritten{\Rightarrow}$1), et \deleted{\unsure{e}}Evidemment 4)$\handwritten{\Rightarrow}$2), d'ailleurs 2)$\handwritten{\Rightarrow}$3) en vertu de 3.5. , grâce au fait que les $\mathrm A_i$ sont noethériens. Reste à prouver 1)$\handwritten{\Rightarrow}$4) et 3)$\handwritten{\Rightarrow}$4).

Supposons 1) vérifié, alors en vertu de 1.7. $u_0^{\handwritten{-}1}:\mathrm C_0\handwritten{\to}\mathrm B_0$ se prolonge en un homomorphisme d'algèbres top $v:\mathrm C\handwritten{\to}\mathrm B$ (grâce aux hypothèses \deleted{\unsure{b}}\handwritten{c}) et \handwritten{d}) ). Or on a déjà noté que les $\mathrm B_i$ sont plats sur les $\mathrm A_i$. Le lemme \deleted{3}\handwritten{4}.2. (i) et (iii) (où on échange les rôles de $\mathrm B,\mathrm C$)
% archive sheet number: 11
